% Article: Conhecimento Comodidy?

\titulo{A Era do Conhecimento}
\documentclass{'../../common-shared/cls/Article_Class_RevInfor'}

%%% Caminho das imagens comuns a todos os artigos
\graphicspath{{../../common-shared/img/}}

%%% Titulo do nosso texto
\titulo{O Conhecimento agora é um comodity?}

\authorRevInfor
{Antonio}
{Wagner}

%%%% Se a gente quizer o banner de cima comum aos artigos da revista infor
%%%% basta descomentar abaixo
%\cabecalhoRevInfor{cabecalho-img-revistainfor.png}

\begin{document}
\textual

\section{Relevância do Tema e o nosso momento histórico}
Será que os cidadãos que viveram na época da Revolução Industrial tinham a consciência do momento histórico e do potencial transformador daquele momento?
Estamos vivendo um momento histórico tal impactante quanto e havendo a oportunidade de refletir sobre os impactos tremendos dessa transformação histórica não é só uma justificativa mas uma responsabilidade histórica também.

\section{Delimitação do estudo}
Dada a relevância e a urgência de estudo do tema, a vontade é de não limitar e sim estudar todas as linhas epistemológicas possíveis para chegar às mais ricas reflexões. Infelizmente isso não é possível dado restrições de tempo etc.

Sendo assim há que se escolher se vamos estudar os impactos economicamente, socialmente, filosoficamente ou que conjunto desses escopos.

Além disso, faremos isso a partir de uma linha epistemológica histórico-crítica, anárquica ou que linha de raciocínio?

Precisamos justificar cada decisão, mas também não descartar que diferentes decisões nesse sentido certamente resultariam em diferentes conclusões e implicações.



\section{O processo histórico da disponibilização das informações}
antes a informação era escassa mas com o advento da imprensa, da internet e da iao isso mudou


\section{A informação e o conhecimento}
O conhecimento, historicamente, sempre esteve em círculos restritos
Se a informação se disseminou podemos dizer o mesmo do conhecimento?

\section{Temas específicos a serem abordados}
Qual a diferença entre o advento da imprensa, da internet e agora da ia?
A avalanche informacionais proveniente da imprensa e internet nos força a pensar na curadoria do que se lê? A IA está fazendo curadoria? Como se posicionar diante disso?



\section{Aplicações}
Aplicações para a sociedade de hoje

\section{Glossário de conceitos}
\begin{enumerate}
  Informação vs Conhecimento
  



